% chktex-file 0
% ===== 第 1 章：早期 Eullossetch 概况 =====
\chapter{早期 Hekov半岛 概况}
现在所谓的\textit{Eullossetch Republika}，是1930年之后的产物。在最早的时期，根据最新的考古发现，大约在十万年前，我们的\textit{Hekov}半岛上，就有早期\textit{Ghoul}族部落、\textit{Fur}族部落，还有一些零零散散的人类部落的活动痕迹。在\textit{Rokadia}河流的下游，靠近\textit{Ignisis}海的区域，考古学家发掘出了大量的砍削器、砍砸器，还有骨制鱼钩、骨矛，证明彼时的诸个部落，以采集和狩猎，以及极小范围的渔业为生。根据目前的主流结论，当时大部分的部落，都住在简易的窝棚里面。距今约四万年前，生产力显著提升，出现了鱼叉，还有原始弓箭等。在如今的\textit{Czelivi}郊外，以及\textit{Lavit}河沿岸，发现了以鲸鱼皮和\textit{Hekov}地区特有的狮鹫羽毛（已经通过相关形态学鉴定）搭建的圆形居住遗址。这说明了两点：其一，当时的部落已经具备迁徙能力，能够从\textit{Ignisis}海岸向内陆迁徙；其二，当时的部落已经有了相对稳定的居住点和抵御严寒的能力。经过一些学者的研究与推断，普遍认为，这些部落已经实现了所谓的\textbf{母系-舅系氏族公社}，奉行\textbf{原始生产资料公有制}，也可能出现了简单的职业分工。在\textit{Lavit}河的沿岸考古基地，我们发现了一根距今三万年、已经愈合的人类大腿骨。这说明当他受伤之后，同部落的人照顾他，让他休息，为他带来食物。在同一考古基地，还发现了已经碳化的研磨草药（也有鉴定观点认为，这是出于食用目的的研磨）。也许，他的同胞们，还会在露水爬上枝头的清晨，为他换药，直到他康复，又能重新行走在大地上。

那是我们的童年时期。有部分心理史学派认为，就像一个人在成长过程中，童年期的经历会影响其成年之后的行为，早期的\textit{Hekov}半岛部落经历了长期的迁徙和生存斗争，让这里的Ghoul族和Fur族在内心深处形成了坚韧的性格，和对生存的执着。但这种说法遭受了大量的质疑，认为这是过度曲解了部落时期文化的影响。再加上后来的各种雇佣兵出口和奴隶贸易带来的换种，这种影响到底有多少，就算有，也是微乎其微、趋近于零。更何况，别的地区的部落，也经历了类似的迁徙和生存斗争，难道只有\textit{Hekov}半岛的部落，才会形成这种性格呢？

童年这个比喻，虽然在修辞学上很合适，但如果用它来解释种族性格的话，显然是不合适的。

\textit{（第一章未完待续）}